% \begin{savequote}
% Alles in Ordnung, auch wenn Sie in Panik geraten sind und sich beschissen fühlen.
% Ich kenne das. Es ist furchtbar, und es wird nicht das letzte Mal sein.
% Aber vergessen Sie nicht, was Sie gesehen haben. Regel Nr. 3: Bei Herzstillstand
% zuerst den \EE{eigenen} Puls fühlen.
% 
% % \begin{flushright}
% \qauthor{Dickie, in: Samuel Shem: \emph{House of God.} \cite{HouseOfGodDe}}
% % \end{flushright}
% 
% \includegraphics[width=4cm]{./images/wikipedia-pd/AED_Symbol}
% \end{savequote}
\chapter
[\CO{[AIR]} Erweitertes Airwaymanagement]
{\CC{[AIR]}\\Erweitertes Airwaymanagement}
% \AasRsN{RS VIII}
\label{AIR}
\label{topic:airway}
\label{topic:airwaymanagement-erweitert}

\ChapterIntro
{
~~~
}{%
\begin{description}
  \item[Maintainer:] Sebastian Gabriel
  \item[Autoren:] Diverse
  \item[Reviewer:] Standardreviewprozess
  \item[Version:] {\DocVersion} ({\DocVersionInfo})
  \item[SHA1:] {\DocGitCommit}
\end{description}

{\TxTeilkapitelAASS}
}

%\SectionUnderConstruction{}
\clearpage
\section{Einleitung}

\Layout{0}{\TxBeschreibung}
{Nach der Kontrolle der Vitalfunktionen kann sich
herausstellen, dass der Patient \EE{kein Bewusstsein} und
\EE{keine normale Atmung} hat, d.\,h. \EE{kein Kreislauf
vorhanden} ist. In dieser Situation liegt die
Notfalldiagnose
\T{Atem-/Kreislaufstillstand} vor. Die einzig
mögliche Therapie ist die (sofortige)
\EE{Herz-Lungen-Wiederbelebung} (\EE{Reanimation}). Dabei
versucht der Helfer den Kreislauf des Patienten \enquote{von
außen} so gut wie möglich aufrecht zu
erhalten. Dies geschieht einerseits durch Beatmung
um dem Körper Sauerstoff zuzuführen und andererseits durch Herzdruckmassage um die
mechanische Herztätigkeit zu übernehmen. Zusätzlich wird versucht mittels
Elektroschocks (Defibrillation) die körpereigene elektrische
Erregung der Herztätigkeit zu \enquote{normalisieren}.
Beatmung, Herzdruckmassage und Defibrillation
bilden die \EE{drei Säulen der Reanimation}.}{
\begin{itemize}
  \item Kein Bewusstsein
  \item Keine normale Atmung
  \item → Kein Kreislauf
  \item → Herz-Lungen-Wiederbelebung (Reanimation)
\end{itemize}
}{}{}{}


\section{Anatomie und Physiologie der Atmung und
Atemwege}
\begin{itemize}
  \item Ventilation, Oxygenation, Respiration;
\end{itemize}
---
\begin{itemize}
   \item die anatomischen Strukturen der Atmung bei Kinder und
  Erwachsene benennen und beschreiben können. 
  \item die
  Atemmechanik verstehen und erklären können.
  
  \item die Physiologie der Atmung verstehen und erklären
  können.
\end{itemize}



\section{Erkrankungen und
Verletzungen der Atemwege und der Atmung}
\begin{itemize}
  \item Krankheitsbilder der Atemwege und der Atmung bei Kinder
  und Erwachsene erkennen und beschreiben können sowie die
  rettungs- und notfallmedizinischen Maßnahmen erklären
  können.
  
  \item Verletzungen der Atemwege und der Atmung bei Kinder und
  Erwachsene erkennen und beschreiben können sowie die
  rettungs- und notfallmedizinischen Maßnahmen erklären
  können.
  
\end{itemize}



\section{Evaluierung der Atemwege 
und der Atmung}
\begin{itemize}
\item die Atemwege und die Ventilation evaluieren und geeignete
  Maßnahmen zur Behebung von Störungen kennen sowie die
  Anwendung dieser Maßnahmen beschreiben können.
\end{itemize}









\section{Ergänzende Sauerstofftherapie bei Erwachsenen und Kindern}
\begin{itemize}
  \item Sauerstoffquellen, Sauerstoffdosierung und Regelung der
  Sauerstoffflusses, Sauerstoffverabreichungsprodukte
  \item und
  Geräte (inkl. Befeuchter) kennen und anwenden können.
\end{itemize}

\subsection{Das Prinzip der Überdruckbeatmung}
\begin{itemize}
  \item Ventilation
  
  \item Mechanik der Überdruckbelüftung \item Magenüberblähung und
  Maßnahmen dagegen
  
  \item Mund-zu-Mund Beatmung,
  
  \item Maskenbeatmung
  
  \item Beatmungsgeräte
  
  \item verschiedene Beatmungsformen und Beatmungskurven
\end{itemize}
---
\begin{itemize}
  \item die Grundlagen und Prinzipien der Unterdruck\item und
  Überdruckbeatmung kennen und beschreiben können.
  
  \item die verschiedenen Möglichkeiten der präklinischen
  Beatmung kennen und anwenden können.
  
  \item an einem Notfalltransportrespirator Grundeinstellungen
  vornehmen können und anwenden können.
   \item die verschiedenen Beatmungsformen kennen und beschreiben
  können.
  \item die Indikation, Kontraindikation sowie 
  mögliche
  Komplikationen bei der Beatmung und Anwendung der
  verschiedenen Techniken kennen und erklären können.
\end{itemize}

\subsubsection{Beatmung mit Beutel und Notfallrespirator}

\section{Atemwegsmanagement}

\subsection{Basistechniken des
Atemwegsmanagement} 



\Layout{2}
{Lagerung}
{}
{}
{}
{}
{}

\Layout{2}
{Absaugen}
{}
{}
{}
{}
{}

\subsubsection{Manuelle Techniken}

\Layout{2}
{Esmarch-Handgriff}
{}
{}
{}
{}
{}

\subsubsection{Atemwegshilfen}

\subsubsection{Sonderfall: Stoma}

\subsection{Endotracheale Intubation}
\subsection{Supraglottische Atemwegssicherungen}
\subsubsection{Larynxtubus}
\subsubsection{Larynxmaske}
\subsubsection{Combitubus}

\subsection{der schwierige Atemweg}

\section{Algorithmen}
\subsection{ASBÖ-Algorithmus Intubation}
---
\begin{itemize}
  \item die Technik der endotrachealen Intubation mit ihren Vor-
  und Nachteilen kennen und beschreiben können.
  
  \item bei der Vorbereitung und Durchführung der Intubation
  einem Arzt oder NFS-NKI assistieren können. \item die
  alternativen Atemwege kennen und bei der Anwendung
  assistieren können.
  
  \item die Indikation, Kontraindikation sowie die Komplikationen
  alternativer Atemwege kennen und beschreiben können.
\end{itemize}


---
\begin{itemize}
  \item die Basistechniken des Atemwegsmanagements beherrschen und
  anwenden können sowie die Indikation, Kontraindikation und
  die Komplikationen beschreiben können.
\end{itemize}




\section{Szenarientraining und Fallbeispiele}

Fallbeispiele: \enquote{Management der Atemwege} -Basismaßnahmen
\begin{itemize}
  \item erweiterte Maßnahmen -alternative Atemwege insbesondere
  die Anwendung von Larynxtuben
\end{itemize}
---

\begin{itemize}
\item im Szenarientraining anhand der gültigen Algorithmen des
  ASBÖ (NKA, NKV, NKI -- Arzneimittellisten) die praktische
  Anwendung der Basismaßnahmen sowie der erweiterten Maßnahmen
  ausgiebig trainieren.
  \item die Assistenz zur Intubation oder alternativer Atemwege
  beherrschen.
\item die fachgerechte Anwendung von Larynxtuben bei Kindern undErwachsene beherrschen. 
\end{itemize}

\section{Verantwortung, kritisches Denken, Sicherheit und Hygiene.}
\begin{itemize}
  \item die Unterschiede, Probleme und Gefahren beim
  Atemwegsmanagement im Schulungsraum (Skilltraining), in der
  Präklinik und im klinischen Setting kennen und erklären
  können.
  
  \item die Maßnahmen und Handfertigkeiten in der Präklinik
  kritisch überdenken (von welchen Maßnahmen profitiert der
  Patient zum gegenwärtigen Zeitpunkt?).
  
  \item alle sicherheitsrelevanten und hygienischen Aspekte in der
  präklinischen Versorgung kennen und wiedergeben können.
  
\end{itemize}
\section{Rechtliche Grundlagen und Einsatzdokumentation}\begin{itemize}
\item die rechtlich Aspekte für seine Tätigkeit als
  Sanitäter (SanG, SanAV und SanAFV) kennen und beschreiben
  können
  \item seine Einsätze vollständig und nachvollziehbar
  dokumentieren können.
\end{itemize}


